\section{Motivation \& Zielsetzung}

\section{Phasen und Ablauf der Studie}

Die Studie gliedert sich ich zwei Phasen: Die Pilotphase und daran anschließend die Hauptstudie. Die Pilotphase gilt der Vorbereitung auf die Hauptstudie und, wie im folgenden Kapitel genauer erläutert, der Abstimmung der einzelnen Elemente. Als Ausgangspunkt wurden Fragebögen, Programmieraufgabe und Bewertungsraster initial unter Einbezug wissenschaftlicher Quellen erstellt und im Zuge des Pilots iterativ im Sinne der agilen Arbeitsweise angepasst.

\subsection{Pilotphase}
Die Pilotphase wird für die Erpobung der einzelnen Teilaspekte der Hauptstudie herangezogen. Besonderes Augenmerk liegt auf der Kalibrierung der Prüfungsaufgaben. Diese sollen so gestaltet sein, dass sie für die Gruppe mit LLM-Unterstützung nicht zu leicht lösbar sind und für die Gruppe ohne LLM-Unterstützung nicht zu überfordernd werden.

In der Pilotstudie wird darauf Acht gegeben, dass (1) die Aufgabe in der vorgegebenen Zeit machbar ist, und (2) die Studierenden ein laut Bewertungsraster mit "bestanden" bewertbares Ergebnis erreichen.

Sollten die Datenerhebungen aus der Pilotstudie auf Probleme in der Aufgabenstellung hinweisen, etwa einen unangemessenen Schwierigkeitsgrad oder eine nicht passende Bearbeitungszeit (zu kurz oder zu lang), können diese vorläufigen Erkenntnisse genutzt werden, um die Aufgaben für die Hauptstudie gezielt anzupassen.

Als Gruppengröße werden zwei Personen als Stichprobe pro Gruppe veranschlagt, somit liegt die Anzahl des Pilots bei \texttt{n=4}. Die Pilot-ProbandInnen sollen einen möglichst ähnlichen Hintergrund wie die Studierenden der Hauptstudie haben und wurden dementsprechend ausgewählt. Als Ablauf des Pilots erfolgte in den folgenden Schritten aus Abbildung \ref{fig:flow-pilotstudie}:

\begin{figure}[htbp]
  \centering
  \begin{tikzpicture}[
    node distance=9mm and 7mm,
    process/.style={
      draw=black!80, fill=white, line width=0.8pt,
      text width=\dimexpr(\linewidth-14mm-32pt)/3\relax,
      minimum height=24mm, inner xsep=5pt, inner ysep=5pt,
      align=center, font=\fontsize{10pt}{14pt}\selectfont
    },
    flow/.style={-{Stealth[length=2.5mm]}, line width=0.8pt}
  ]
    \node[process] (questionnaire1) {Bearbeitung Fragebogen 1:\\Selbsteinschätzung};
    \node[process, right=of questionnaire1] (environment) {Einführung in die\\Entwicklungsumgebung};
    \node[process, right=of environment] (presentation) {Präsentation der Aufgabe\\und Beantwortung\\anfallender Fragen};
    \node[process, below=of presentation] (task) {Bearbeitung der Aufgabe\\durch die TeilnehmerInnen};
    \node[process, left=of task] (questionnaire2) {Bearbeitung Fragebogen 2:\\UEQ nach der Aufgabe};
    \node[process, left=of questionnaire2] (collection) {Screen Recording\\speichern und\\finalen Code einsammeln};
    \draw[flow] (questionnaire1) -- (environment);
    \draw[flow] (environment) -- (presentation);
    \draw[flow] (presentation) -- (task);
    \draw[flow] (task) -- (questionnaire2);
    \draw[flow] (questionnaire2) -- (collection);
  \end{tikzpicture}
  \caption{Ablauf der Pilotstudie}
  \label{fig:flow-pilotstudie}
\end{figure}

Nach der Vorstellung erfolgt als erstes die Bearbeitung des ersten Fragebogens, bei dem es um die Selbsteinschätzung der Studierenden geht. Danach werden die Teilnehmenden in die Entwicklungsumgebung VS Code eingeführt und sichergegangen, ob alle Vorraussetzungen vorhanden sind: Python und Intellisense sind installiert und die Bildschirmaufzeichnung ist startklar. Anschließend wird die Aufgabe vorgestellt und erklärt wie die lokalen Tests ausgeführt werden können. Nach der Bearbeitung aller Teilaufgaben folgt der zweite Fragebogen. Dieser behandelt das Thema der User Experience. Der letzte Schritt stellt das Einsammeln des Codes und der Bildschirmaufnahme dar.

Während der Druchführung des Pilots werden von der Studienleitung verbesserungswürdige Themen niedergeschrieben, um diese in agiler Art und Weise vor erneuter Durchführung einzuarbeiten. Die Auswertung der Pilotphase wird in Abschnitt \ref{cha:pilot_auswertung} beschrieben.

\subsection{Hauptstudie}

Der Ablauf der Hauptstudie wurde in weiten Teilen dem des Pilots in Abbildung \ref{fig:flow-pilotstudie} nachempfunden. Es wurde für die Durchführung eine granulare Checkliste in chronologischer Reihenfolge erstellt, die in Anhang~\ref{app:checkliste-hauptstudie} enthalten ist. Hier können explizit die Schritte der Studie eingesehen werden. Im folgenden Abschnitt soll auf Aspekte eingegangen werden, die Änderungen oder Neuerungen im Ablauf der Hauptstudie nach sich zogen.

Neu eingeführt in der Hauptstudie wurde die Nummernvergabe am Anfang zur Pseudonymisierung der erhobenen Daten. Die Studierenden waren angehalten, als ersten Programmpunkt eine Nummer zu ziehen, mit der in Folge alle Abgaben und Fragebögen zu kennzeichnen waren. Für die Zuordnung wurde eine Tabelle verwendet, deren Vorlage in Anhang~\ref{app:zuordnungstabelle} enthalten ist. Hier wurden Geburtsdatum, Gender sowie Vor- und Nachname der Teilnehmenden erfasst. Die Auswertung soll unter Verwendung der vergebenen Nummern erfolgen, sodass für Außenstehende kein Rückschluss auf die tatsächliche Person möglich ist. Mithilfe der Zuordnungstabelle kann jedoch immer eine Zuordnung der Daten zu den jeweiligen Teilnehmenden durchgeführt werden. Diese Tabelle wird nicht im Zuge der Arbeit veröffentlicht.

Als weiteres Learing aus dem Pilot wurde das Erstellen eines Github Accounts oder der Login in denselben als Programmpunkt ergänzt. Damit verbunden war auch das Repository auf Github zu forken und den den Github codespace zu starten.

Ein Punkt, der im Zuge der Pilotphase an Wichtigkeit erlangte und sich somit als eigener Programmpunkt sich wiederfand war die Sachgemäße und einheitliche Einrichtung der Hilfsmittel. Für die KI-Gruppe war es die Gratis-Version von ChatGPT, ohne angemeldet zu sein. Für die Websuche-Gruppe wurde DuckDuckGo verwendet, da diese Suchmaschine die Deaktivierung von AI-Features bei einer Websuche ermöglicht. Letzteres wurde im Zuge der Pilotphase angepasst, wie in \ref{cha:Adaptionen} erläutert, und anschließend als Programmpunkt in die Hauptstudie übernommen.

Um eine reibungslose Durchführung der einzelnen Phasen der Hauptstudie für die Studierenden zu gewährleisten wurde eine Padlet-Leinwand mit Links und Notizen bereitgestellt. Die in Anhang~\ref{app:padlet} dargestellte Leinwand beinhaltet die Links zum persönlichen Fragebogen, zum GitHub-Repository des \gs und zum abschließenden Fragebogen zur User Experience. Zusätzlich enthält sie den Hinweis, die Änderungen im Codespace zu committen und die persönliche Nummer in der Commit-Nachricht anzugeben. Für die Bildschirmaufzeichnung wird Dateinamensschema sowie ein Link zum Hochladen in einen Google-Drive-Ordner bereitgestellt.

Die Aufgaben, die von den Studierenden im Zuge der Studie zu erledigen sind, wurde ebenfalls angepasst und erweitert. Im Ablauf der Studie gab es deshalb jedoch keine Änderung, die Erledigung der Aufgaben blieb schematisch an dem Punkt wie in \ref{fig:flow-pilotstudie} dargestellt.
